\chapter{Calcul d'Itô}\label{ch:qm:ito-calculus}

Un jeune chercheur soumet à un backtest une règle de suivi de tendance sur dix ans de cours de clôture quotidiens~: détenir une
position proportionnelle à l'écart du cours de clôture au-dessus de sa moyenne mobile sur 20 jours. Le
ratio de Sharpe vaut 5.4. La règle lit le cours de clôture, et le backtest lui impute le
rendement qui entre dans ce même cours~; décalé d'un jour, de sorte qu'une position décidée à un cours
rapporte le rendement du lendemain, le ratio de Sharpe vaut $-0.04$. Les «~cours~» étaient une pure
marche aléatoire, sans rien à trouver. Le premier backtest n'était pas une \omterm{def:qm:ito-calculus:integral}{intégrale d'Itô}~: son intégrande
examinait l'accroissement par lequel elle était multipliée, et collectait la covariation de la position avec le
cours, une quantité qui existe sur toute trajectoire aléatoire. Ce chapitre construit l'\omterm{def:qm:ito-calculus:integral}{intégrale d'Itô},
dont l'intégrande est fixée avant l'arrivée de l'accroissement, démontre la formule d'Itô, la règle de dérivation des fonctions composées,
avec un terme du second ordre que le reste de la série utilise à chaque page, et se termine par les
deux théorèmes structurels (la caractérisation de Lévy et la représentation par \omterm{def:qm:probability-at-speed:martingale}{martingales}) qui
sous-tendent la tarification du livre~5.

\section{L'intégrale d'Itô}

Une stratégie de trading détient $H_t$ unités d'un actif dont le cours suit un \omterm{def:qm:brownian-motion:bm}{mouvement brownien}~; ses
gains sur $[0, T]$ devraient être $\int_0^T H_t\,dW_t$. Les trajectoires de $W$ sont à variation totale infinie
(le \cref{thm:qm:brownian-motion:qv}), donc l'intégrale ne peut pas être définie trajectoire par trajectoire comme une
intégrale de Stieltjes, et le choix du point d'évaluation dans chaque intervalle compte à la limite.
L'\omterm{def:qm:ito-calculus:integral}{intégrale d'Itô} évalue $H$ au point de gauche, seul choix qu'un trader puisse faire.

\begin{definition}[Processus élémentaire, intégrale d'Itô]\label{def:qm:ito-calculus:integral}
Un \emph{processus élémentaire}\index{processus élémentaire} est $H_t = \sum_{k=0}^{n-1}h_k\mathbf 1_{(t_k,
t_{k+1}]}(t)$ avec $0 = t_0 < \dots < t_n = T$ et chaque $h_k$ borné et $\mathcal
F_{t_k}$-mesurable. Son \emph{intégrale d'Itô}\index{intégrale d'Itô} est $\int_0^t H_s\,dW_s = \sum_k
h_k(W_{t_{k+1}\wedge t} - W_{t_k\wedge t})$. Pour un $H$ adapté avec $\E\int_0^T H_s^2\,ds <
\infty$, l'intégrale d'Itô est la limite dans $L^2$ des intégrales des processus élémentaires $H^{(n)}$ avec
$\E\int_0^T(H^{(n)}_s - H_s)^2\,ds \to 0$.
\end{definition}

\begin{theorem}[Isométrie d'Itô]\label{thm:qm:ito-calculus:isometry}
Pour $H$ adapté avec $\E\int_0^T H_s^2\,ds < \infty$, $I_t = \int_0^t H_s\,dW_s$ est une
\omterm{def:qm:probability-at-speed:martingale}{martingale} continue de carré intégrable, avec $I_0 = 0$ et
\[
\E\Bigl[\Bigl(\int_0^T H_s\,dW_s\Bigr)^2\Bigr] = \E\int_0^T H_s^2\,ds, \qquad [I]_t = \int_0^t H_s^2\,ds.
\]
\end{theorem}

\begin{proof}
Pour un \omterm{def:qm:ito-calculus:integral}{processus élémentaire}, avec $\Delta_k = W_{t_{k+1}} - W_{t_k}$~: les termes croisés $\E[h_jh_k\Delta_j\Delta_k]$,
$j < k$, s'annulent car $h_jh_k\Delta_j$ est $\mathcal F_{t_k}$-mesurable et $\E_{t_k}[\Delta_k] = 0$~;
les termes carrés donnent $\E[h_k^2]\E[\Delta_k^2] = \E[h_k^2](t_{k+1} - t_k)$. Le même calcul sur
$[s, t]$ donne la propriété de \omterm{def:qm:probability-at-speed:martingale}{martingale}. L'intégrale est donc une isométrie des \omterm{def:qm:ito-calculus:integral}{processus élémentaires}
vers $L^2$, se prolonge à leur adhérence, et la propriété de \omterm{def:qm:probability-at-speed:martingale}{martingale} passe aux
limites dans $L^2$ (l'inégalité de Doob, la \cref{prop:qm:probability-at-speed:doob}, fournit une
version continue).
\end{proof}

\begin{example}[L'intégrale de $W$ contre elle-même]\label{ex:qm:ito-calculus:wdw}
Sur une subdivision de $[0, T]$, $\sum_k W_{t_k}\Delta_k = \tfrac12\sum_k(W_{t_{k+1}}^2 - W_{t_k}^2) -
\tfrac12\sum_k\Delta_k^2$. La première somme se télescope en $\tfrac12 W_T^2$ et la seconde tend vers
$\tfrac12 T$ d'après le \cref{thm:qm:brownian-motion:qv}~:
\[
\int_0^T W_t\,dW_t = \tfrac12\bigl(W_T^2 - T\bigr).
\]
Les sommes au point de droite tendent vers $\tfrac12(W_T^2 + T)$, et les sommes au point milieu vers $\tfrac12W_T^2$, la
réponse du calcul ordinaire (la \cref{fig:qm:ito-calculus:sums}).
\end{example}

\begin{omfigure}
\begin{tikzpicture}
\begin{axis}[width=11cm,height=5cm,xmode=log,log basis x=2,
  xlabel={points dans $[0,1]$},ylabel={somme},
  tick label style={font=\small},label style={font=\small},
  xmin=3,xmax=90000,ymin=0.5,ymax=3.6,grid=major,grid style={gray!25},
  legend columns=3,legend style={font=\footnotesize,at={(0.5,-0.3)},anchor=north,draw=none,fill=none,/tikz/every even column/.append style={column sep=8pt}},
  table/col sep=comma]
\addplot[omThm,thick,mark=*,mark size=1.2pt] table[x=n,y=right]{figdata/methods/03-ito-calculus/sums.csv};
\addplot[omProp,thick,mark=square*,mark size=1.2pt] table[x=n,y=mid]{figdata/methods/03-ito-calculus/sums.csv};
\addplot[omDef,thick,mark=triangle*,mark size=1.6pt] table[x=n,y=left]{figdata/methods/03-ito-calculus/sums.csv};
\addplot[gray,dashed,forget plot] table[x=n,y=ito]{figdata/methods/03-ito-calculus/sums.csv};
\addplot[gray,dashed,forget plot] table[x=n,y=rlim]{figdata/methods/03-ito-calculus/sums.csv};
\legend{point de droite,point milieu (Stratonovich),point de gauche (Itô)}
\end{axis}
\end{tikzpicture}

\omcaption{Sommes approchant $\int_0^1 W\,dW$ sur une trajectoire ($W_1 = 2.04$) échantillonnée en $2^k$ points.
Les sommes au point de gauche convergent vers $\tfrac12(W_1^2 - 1) = 1.58$, les sommes au point de droite vers $2.58$
(traits pointillés)~; les sommes au point milieu valent exactement $\tfrac12W_1^2 = 2.08$, quelle que soit la grille. Données~:
le tutoriel du chapitre, à graine fixée.}\label{fig:qm:ito-calculus:sums}
\end{omfigure}

\begin{definition}[Intégrale de Stratonovich]\label{def:qm:ito-calculus:strat}
L'\emph{intégrale de Stratonovich}\index{intégrale de Stratonovich} $\int_0^T X_t\circ dW_t$ est la limite de
les sommes $\sum_k\tfrac12(X_{t_k} + X_{t_{k+1}})\Delta_k$. Pour une semimartingale continue $X$, elle
vaut $\int_0^T X_t\,dW_t + \tfrac12[X, W]_T$.
\end{definition}

L'\omterm{def:qm:ito-calculus:strat}{intégrale de Stratonovich} obéit à la règle de dérivation ordinaire, ce qui explique que les physiciens modélisant un
bruit lisse par un bruit blanc l'utilisent. Elle moyenne l'intégrande sur l'intervalle, ce qu'un
trader ne peut pas faire~: la position est fixée avant le mouvement du cours. La finance utilise l'\omterm{def:qm:ito-calculus:integral}{intégrale d'Itô}.

\section{L'intégrale dans un backtest}

Un backtest sur une grille $t_0 < t_1 < \dots$, avec des positions $\theta_{t_i}$ et des cours $S_{t_i}$, calcule
gains $\sum_i\theta_{t_i}(S_{t_{i+1}} - S_{t_i})$~: la somme au point de gauche, une \omterm{def:qm:ito-calculus:integral}{intégrale d'Itô}, lorsque
$\theta_{t_i}$ n'utilise que l'information disponible en $t_i$. Un backtest qui impute à la position
calculée en $t_{i+1}$ le mouvement entrant dans $t_{i+1}$ calcule à la place la somme au point de droite, et
\[
\sum_i\theta_{t_{i+1}}\Delta S_i - \sum_i\theta_{t_i}\Delta S_i = \sum_i\Delta\theta_i\,\Delta S_i
\longrightarrow [\theta, S]_T:
\]
le gain d'anticipation est la covariation de la position avec le cours, qui n'est pas nulle
sur une marche aléatoire dès que la position réagit au dernier cours. La règle du crochet,
$\theta_t = (S_t - \bar S^{(20)}_t)/c$ avec $\bar S^{(20)}$ la moyenne mobile sur 20 jours et $c$ choisi
de sorte que $\theta$ soit de variance unitaire, et accorde un poids $(L - 1)/L = 0.95$ au dernier accroissement~:
la \cref{fig:qm:ito-calculus:timing} montre le calage temporel et la \cref{fig:qm:ito-calculus:backtest} le
résultat.

\begin{omfigure}
\begin{tikzpicture}[x=3.4cm,y=1cm,font=\small]
\draw[->,thick,gray] (0,0) -- (3.3,0) node[right,black] {temps};
\foreach \x/\l in {0/{clôture $t-1$},1/{clôture $t$},2/{clôture $t+1$}} {\draw[thick] (\x,-0.12) -- (\x,0.12); \node[below=3pt] at (\x,-0.1) {\l};}
\draw[omDef,very thick,-{Stealth}] (1,0.35) to[bend left=25] node[above,align=center] {Itô~: $\theta_t$ gagne\\$S_{t+1} - S_t$} (2,0.35);
\draw[omThm,very thick,-{Stealth},dashed] (1,-0.75) to[bend left=25] node[below=2pt,align=center] {anticipation~: $\theta_t$ gagne\\$S_t - S_{t-1}$} (0,-0.75);
\node[draw=omDef,rounded corners=2pt,fill=omDef!6,inner sep=3pt,font=\footnotesize,align=center] (box) at (0.35,1.55) {position $\theta_t$ calculée\\à partir du cours de clôture $S_t$};
\draw[gray,->] (box.east) -| (1,0.18);
\end{tikzpicture}

\omcaption{Quel mouvement de cours une position peut gagner. Une position calculée à partir du cours de clôture en $t$
gagne le mouvement de $t$ à $t + 1$ (trait plein)~: la somme au point de gauche. Le créditer du mouvement entrant dans
$t$ (pointillés) multiplie la position par un accroissement qu'elle a déjà vu.}\label{fig:qm:ito-calculus:timing}
\end{omfigure}

\begin{proposition}[Le ratio de Sharpe avec anticipation sur une marche aléatoire]\label{prop:qm:ito-calculus:lookahead}
Soit $S$ une marche aléatoire d'accroissements quotidiens $\sigma Z_t$, les $Z_t$ étant des normales centrées réduites indépendantes,
et $\theta_t$ une fonction linéaire des accroissements passés, de variance unitaire et de corrélation $a$ avec le
dernier. Alors les gains de même barre $\theta_t\,\Delta S_t$ ont pour moyenne $a\sigma$ et pour
écart type $\sigma\sqrt{1 + a^2}$, soit un ratio de Sharpe quotidien de $a/\sqrt{1 + a^2}$, tandis que les gains de la barre suivante
ont une moyenne nulle. Pour la règle de moyenne mobile, $a = \frac{(L-1)/L}{(\sum_{m=1}^{L-1}(m/L)^2)^{1/2}}$.
\end{proposition}

\begin{proof}
Écrivons $\theta_t = aZ_t + \sqrt{1 - a^2}\,Y$ avec $Y$ normale centrée réduite et indépendante de $Z_t$. Alors
$\E[\theta_tZ_t] = a$ et $\E[\theta_t^2Z_t^2] = a^2\E[Z^4] + (1 - a^2) = 1 + 2a^2$, donc la variance vaut
$1 + a^2$ (en unités de $\sigma^2$). Pour la règle, $S_t - \bar S_t = \sum_{i=0}^{L-2}\frac{L-1-i}{L}\Delta
S_{t-i}$, dont la variance est $\sigma^2\sum_{m=1}^{L-1}(m/L)^2$ et dont le coefficient sur $\Delta S_t$ est
$(L-1)/L$. Les gains de la barre suivante $\theta_{t-1}\Delta S_t$ ont une moyenne nulle, car $\theta_{t-1}$ est indépendant
de $\Delta S_t$.
\end{proof}

Pour $L = 20$, $a = 0.382$ et le ratio de Sharpe annualisé de même barre vaut $\sqrt{252}\times 0.382/
\sqrt{1.146} = 5.67$~; sur 400 marches aléatoires simulées de dix ans, il vaut en moyenne 5.69, et le backtest honnête
donne 0.00 avec un écart type de 0.32, le bruit qu'on attend de dix ans d'une
stratégie sans edge.

\begin{omfigure}
\begin{tikzpicture}
\begin{axis}[width=11cm,height=5cm,
  xlabel={années},ylabel={P\&L cumulé (\% du notionnel unitaire)},
  tick label style={font=\small},label style={font=\small},
  xmin=0,xmax=10,ymin=-100,ymax=1050,grid=major,grid style={gray!25},
  legend columns=3,legend style={font=\footnotesize,at={(0.5,-0.3)},anchor=north,draw=none,fill=none,/tikz/every even column/.append style={column sep=8pt}},
  table/col sep=comma]
\addplot[omThm,very thick] table[x=year,y=cheat]{figdata/methods/03-ito-calculus/backtest.csv};
\addplot[omProp,thick,dashed] table[x=year,y=cov]{figdata/methods/03-ito-calculus/backtest.csv};
\addplot[omDef,very thick] table[x=year,y=honest]{figdata/methods/03-ito-calculus/backtest.csv};
\legend{même barre (anticipation),covariation $[\theta{,}S]$,barre suivante (Itô)}
\end{axis}
\end{tikzpicture}

\omcaption{La règle de tendance à 20 jours sur une marche aléatoire de dix ans avec une volatilité quotidienne de 1~\%. Créditée
de la même barre, elle rapporte 958~\% du notionnel (ratio de Sharpe 5.4), presque exactement la covariation de
la position et du cours (965~\%)~; créditée de la barre suivante, elle rapporte $-7\%$ (ratio de Sharpe $-0.04$).
Données~: le tutoriel du chapitre, à graine fixée.}\label{fig:qm:ito-calculus:backtest}
\end{omfigure}

\section{Formule d'Itô en une et plusieurs dimensions}

\begin{definition}[Processus d'Itô, covariation quadratique]\label{def:qm:ito-calculus:process}
Un \emph{processus d'Itô}\index{processus d'Itô} est $X_t = X_0 + \int_0^t\mu_s\,ds + \int_0^t\sigma_s\,dW_s$, qui s'écrit
$dX_t = \mu_t\,dt + \sigma_t\,dW_t$, avec $\mu$ et $\sigma$ adaptés tels que $\int_0^T|\mu_s|\,ds$ et
$\int_0^T\sigma_s^2\,ds$ soient finies presque sûrement. La \emph{covariation quadratique}\index{covariation quadratique}
de deux processus continus est la limite en probabilité $[X, Y]_t =
\lim\sum_k(X_{t_{k+1}} - X_{t_k})(Y_{t_{k+1}} - Y_{t_k})$~; pour des processus d'Itô pilotés par $W^1, W^2$
avec $d\langle W^1, W^2\rangle_t = \rho\,dt$, $d[X, Y]_t = \sigma^X_t\sigma^Y_t\rho\,dt$.
\end{definition}

Les règles de calcul sont $dt\,dt = 0$, $dt\,dW = 0$, $dW^i\,dW^j = \rho_{ij}\,dt$~: seuls les produits de
deux accroissements browniens subsistent au premier ordre en $dt$.

\begin{theorem}[Formule d'Itô]\label{thm:qm:ito-calculus:ito}
Si $X$ est un \omterm{def:qm:ito-calculus:process}{processus d'Itô} et $f \in C^{1,2}([0,T]\times\R)$, alors
\[
df(t, X_t) = \partial_tf\,dt + \partial_xf\,dX_t + \tfrac12\partial_{xx}f\,d[X]_t
= \bigl(\partial_tf + \mu_t\partial_xf + \tfrac12\sigma_t^2\partial_{xx}f\bigr)dt + \sigma_t\partial_xf\,dW_t.
\]
Pour $X = (X^1, \dots, X^d)$ et $f \in C^{1,2}$, $df = \partial_tf\,dt + \sum_i\partial_if\,dX^i +
\tfrac12\sum_{i,j}\partial_{ij}f\,d[X^i, X^j]$~; en particulier $d(XY) = X\,dY + Y\,dX + d[X, Y]$.
\end{theorem}

\begin{proof}[Démonstration partielle]
Pour $f(x)$ à dérivées bornées et $\mu, \sigma$ simples, la formule de Taylor sur une subdivision donne
$f(X_T) - f(X_0) = \sum f'(X_{t_k})\Delta X_k + \tfrac12\sum f^{\prime\prime}(\xi_k)(\Delta X_k)^2$. La première somme
converge vers $\int f'(X)\,dX$. Dans la seconde, $f^{\prime\prime}(\xi_k)$ peut être remplacé par $f^{\prime\prime}(X_{t_k})$ au
coût borné par le module de continuité de $f^{\prime\prime}$ multiplié par $\sum(\Delta X_k)^2$, qui tend vers zéro,
et $\sum f^{\prime\prime}(X_{t_k})((\Delta X_k)^2 - \sigma_{t_k}^2\Delta t_k)$ tend vers zéro dans $L^2$ comme dans
le \cref{thm:qm:brownian-motion:qv}. Le cas général de $f$, $\mu$, $\sigma$ découle par localisation et
approximation (Karatzas et Shreve, 1991, §3.3).
\end{proof}

\begin{example}[Le logarithme d'un mouvement brownien géométrique]\label{ex:qm:ito-calculus:log}
Pour $dS = \mu S\,dt + \sigma S\,dW$, avec $f = \ln$~: $d\ln S = dS/S - \tfrac12 d[S]/S^2 = (\mu -
\tfrac12\sigma^2)dt + \sigma\,dW$. Le rendement espéré vaut $\mu$, mais le rendement logarithmique espéré, soit le taux
auquel la médiane croît, vaut $\mu - \tfrac12\sigma^2$. Avec $\mu = 10\%$ et $\sigma = 40\%$, la
moyenne de $S_{10}/S_0$ vaut $e^{1} = 2.72$ et sa médiane $e^{0.2} = 1.22$
(la \cref{fig:qm:ito-calculus:gbm})~; l'érosion par la volatilité des fonds à effet de levier (livre~1,
chapitre~14) est la même correction appliquée deux fois.
\end{example}

\begin{omfigure}
\begin{tikzpicture}
\begin{axis}[width=11cm,height=5cm,ymode=log,
  xlabel={années},ylabel={$S_t/S_0$},
  tick label style={font=\small},label style={font=\small},
  xmin=0,xmax=10,ymin=0.2,ymax=8,grid=major,grid style={gray!25},
  ytick={0.25,0.5,1,2,4,8},log ticks with fixed point,
  legend columns=3,legend style={font=\footnotesize,at={(0.5,-0.3)},anchor=north,draw=none,fill=none,/tikz/every even column/.append style={column sep=8pt}},
  table/col sep=comma]
\addplot[name path=hi,gray!40,forget plot] table[x=t,y=p90]{figdata/methods/03-ito-calculus/gbm.csv};
\addplot[name path=lo,gray!40,forget plot] table[x=t,y=p10]{figdata/methods/03-ito-calculus/gbm.csv};
\addplot[gray!15,forget plot] fill between[of=hi and lo];
\addplot[omThm,very thick] table[x=t,y=mean_th]{figdata/methods/03-ito-calculus/gbm.csv};
\addplot[omDef,very thick] table[x=t,y=median_th]{figdata/methods/03-ito-calculus/gbm.csv};
\addplot[black,only marks,mark=*,mark size=1.3pt] table[x=t,y=median]{figdata/methods/03-ito-calculus/gbm.csv};
\legend{moyenne $e^{\mu t}$,médiane $e^{(\mu - \sigma^2/2)t}$,médiane simulée}
\end{axis}
\end{tikzpicture}

\omcaption{\omterm{def:qm:brownian-motion:bm}{Mouvement brownien} géométrique avec $\mu = 10\%$ et $\sigma = 40\%$~: la moyenne croît de 10~\% par
an, la médiane de 2~\%, et après dix ans le 10\up{e} centile (bande grisée~: du 10\up{e} au 90\up{e}) vaut
0.24. Données~: 100\,000 trajectoires simulées exactement, à graine fixée.}\label{fig:qm:ito-calculus:gbm}
\end{omfigure}

\section{Martingales locales et exponentielle stochastique}

Une \omterm{def:qm:ito-calculus:integral}{intégrale d'Itô} $\int H\,dW$ pour laquelle on sait seulement que $\int_0^TH_s^2\,ds < \infty$ presque
sûrement n'est pas nécessairement intégrable, encore moins une \omterm{def:qm:probability-at-speed:martingale}{martingale}~; elle en est une localement, pour une suite de \omterm{def:qm:probability-at-speed:stopping}{temps d'arrêt} convenable.

\begin{definition}[Martingale locale]\label{def:qm:ito-calculus:local}
Un \omterm{def:qm:probability-at-speed:filtration}{processus adapté} $M$ est une \emph{martingale locale}\index{martingale locale} s'il existe des \omterm{def:qm:probability-at-speed:stopping}{temps
d'arrêt} $\tau_n \uparrow \infty$ tels que chaque processus arrêté $M_{t\wedge\tau_n}$ soit une \omterm{def:qm:probability-at-speed:martingale}{martingale}.
\end{definition}

\begin{proposition}[Les martingales locales positives sont des surmartingales]\label{prop:qm:ito-calculus:super}
Une \omterm{def:qm:ito-calculus:local}{martingale locale} minorée, en particulier positive, est une \omterm{def:qm:probability-at-speed:martingale}{surmartingale}. Une
\omterm{def:qm:ito-calculus:local}{martingale locale} $M$ telle que $\E[\sup_{s\le t}|M_s|] < \infty$ pour tout $t$ est une \omterm{def:qm:probability-at-speed:martingale}{martingale}.
\end{proposition}

\begin{proof}
Si $M \ge 0$, le lemme de Fatou donne $\E_s[M_t] = \E_s[\lim_n M_{t\wedge\tau_n}] \le \liminf_n
M_{s\wedge\tau_n} = M_s$. Sous la condition de domination, le théorème de convergence dominée remplace l'inégalité par une
égalité.
\end{proof}

L'inégalité peut être stricte~: $1/|B_t|$ pour un \omterm{def:qm:brownian-motion:bm}{mouvement brownien} en dimension trois issu d'un point
distinct de l'origine est une \omterm{def:qm:ito-calculus:local}{martingale locale} positive dont l'espérance décroît, soit une
\omterm{def:qm:ito-calculus:local}{martingale locale} \emph{stricte}. Ces processus modélisent des bulles de prix et expliquent pourquoi le chapitre~5 vérifie
qu'un \omterm{def:qm:probability-at-speed:change}{processus densité} candidat est une vraie \omterm{def:qm:probability-at-speed:martingale}{martingale} avant de s'en servir pour opérer un \omterm{def:qm:probability-at-speed:change}{changement de probabilité}.

\begin{definition}[Exponentielle stochastique]\label{def:qm:ito-calculus:exp}
L'\emph{exponentielle stochastique}\index{exponentielle stochastique} d'une semimartingale continue $X$
avec $X_0 = 0$ est $\mathcal E(X)_t = \exp(X_t - \tfrac12[X]_t)$, l'unique solution de $dZ = Z\,dX$,
$Z_0 = 1$.
\end{definition}

D'après la formule d'Itô, $d\exp(X - \tfrac12[X]) = \exp(\cdot)(dX - \tfrac12 d[X] + \tfrac12 d[X]) = Z\,dX$. Avec
$X = \int\gamma\,dW$, c'est la \omterm{def:qm:ito-calculus:local}{martingale locale} positive $\exp(\int\gamma\,dW - \tfrac12\int\gamma^2dt)$,
le \omterm{def:qm:probability-at-speed:change}{processus densité} de tout \omterm{def:qm:probability-at-speed:change}{changement de probabilité} étudié au chapitre~5.

\section{Le théorème de représentation}

\begin{theorem}[Caractérisation de Lévy]\label{thm:qm:ito-calculus:levy}
Une \omterm{def:qm:ito-calculus:local}{martingale locale} continue $M$ telle que $M_0 = 0$ et $[M]_t = t$ est un \omterm{def:qm:brownian-motion:bm}{mouvement brownien}.
\end{theorem}

\admitted

\begin{theorem}[Représentation des martingales]\label{thm:qm:ito-calculus:representation}
Soit $\mathbb F$ la \omterm{def:qm:probability-at-speed:filtration}{filtration} engendrée par un \omterm{def:qm:brownian-motion:bm}{mouvement brownien} $W$. Toute
$\mathbb F$-martingale de carré intégrable $M$ s'écrit $M_t = M_0 + \int_0^tH_s\,dW_s$, pour un unique $H$ adapté vérifiant
$\E\int_0^TH_s^2\,ds < \infty$.
\end{theorem}

\admitted

Les deux résultats sont démontrés dans Karatzas et Shreve (1991, §3.3 et §3.4). Le théorème de représentation
est le fondement mathématique de la réplication~: si le seul aléa est $W$, l'\omterm{def:qm:probability-at-speed:condexp}{espérance conditionnelle} de chaque paiement
est une intégrale stochastique par rapport à $W$, et $H$ est la position qui le reproduit. Le
livre~5, chapitre~1, en tire les notions de portefeuille répliquant et de marché complet~; il
échoue dès qu'il y a des sauts (chapitre~6) ou plus de sources de risque que d'actifs négociés.

\section{Tutoriel~: sommes d'Itô et backtest honnête}\label{tut:qm:ito-calculus:backtest}

\begin{tutorial}
\textbf{Objectif.} Voir en chiffres le point gauche de l'\omterm{def:qm:ito-calculus:integral}{intégrale d'Itô}, vérifier la formule d'Itô sur une
trajectoire, et mesurer le gain lié à l'anticipation d'un backtest sur la même barre. \textbf{État final~:}
les \cref{fig:qm:ito-calculus:sums,fig:qm:ito-calculus:backtest,fig:qm:ito-calculus:gbm} ainsi que les
ratios de Sharpe 5.4 et $-0.04$.
\begin{tutsteps}
\item \textbf{Les intégrales.} Les gains du projet fil rouge sont des sommes à point gauche~; la version
sur la même barre et la covariation qui les sépare tiennent chacune en une ligne.
\omcode{code/firm/stochint/firm_stochint.py}{24}{42}{Gains adaptés, gains sur la même barre et covariation entre eux.}\label{lst:qm:ito-calculus:gains}
\item \textbf{La règle}~: l'écart du cours de clôture à sa moyenne mobile, ramené à une variance unitaire
sous une marche aléatoire, et l'espérance du gain d'anticipation énoncée à la \cref{prop:qm:ito-calculus:lookahead}.
\omcode{code/methods/03-ito-calculus/python/qm_ito.py}{33}{47}{Le signal de tendance et le gain d'anticipation espéré.}\label{lst:qm:ito-calculus:signal}
\item \textbf{Exécuter} \texttt{sums\_table()}, \texttt{ito\_check()} (sur une trajectoire de \omterm{def:qm:brownian-motion:bm}{mouvement brownien} géométrique avec $2^{14}$ pas,
$\ln S_1 = 0.70040$ contre $0.70039$ selon la formule d'Itô et $0.78179$ sans le $-\tfrac12\,d[S]/S^2$
terme), \texttt{backtest()} et \texttt{fig\_ito.py}.
\end{tutsteps}
\textbf{Pour aller plus loin.} Remplacer la règle de tendance par une règle de retour à la moyenne et prédire le signe
du biais d'anticipation~; ajouter un rebond bid-ask d'un tick aux cours et observer ce que le backtest honnête
d'une règle de retour à la moyenne rapporte alors.
\end{tutorial}

\section{Construction~: intégrales stochastiques discrètes}\label{bld:qm:ito-calculus:stochint}

\begin{build}
\bfield{Rôle} Chaque backtest de l'entreprise miniature calcule son P\&L au moyen de ces fonctions,
de sorte qu'une position ne puisse rapporter que le mouvement de prix postérieur à sa décision, et qu'une
anticipation laisse une trace mesurable.
\bfield{Interface} \texttt{gains(theta, s)}~; \texttt{same\_bar\_gains(theta, s)}~;
\texttt{covariation(x, y)}~; \texttt{lookahead\_test(theta, s)} renvoyant l'excédent sur la même barre, la
covariation qu'il égale, une statistique $t$ et l'erreur d'identité~; \texttt{riemann\_sums(h, x,
point)}.
\bfield{Règles} Positions et cours sont alignés sur les mêmes horodatages, et c'est la fonction, non
l'appelant, qui les décale~; l'identité sur la même barre $-$ adaptée $=$ covariation tient à l'arrondi près.
\bfield{Tests d'acceptation} \texttt{code/firm/stochint/tests/}~: sur des marches aléatoires, les gains adaptés de
toute règle sont de moyenne nulle~; les gains sur la même barre de la règle de tendance ont pour moyenne celle de la
\cref{prop:qm:ito-calculus:lookahead}~; l'identité tient à $10^{-12}$ près~; les trois sommes de Riemann de
$\int W\,dW$ convergent vers leurs limites.
\bfield{Pour aller plus loin} Horodatages avec latences (une position devient active quelques microsecondes après son
signal)~; coûts de transaction en fonction de $|\Delta\theta|$.
\end{build}

\begin{omsources}
\item K.~Itô, «~Stochastic integral~», \emph{Proceedings of the Imperial Academy, Tokyo} 20, 1944.
\item K.~Itô, «~On a formula concerning stochastic differentials~», \emph{Nagoya Mathematical
Journal} 3, 1951.
\item R.~L.~Stratonovich, «~A new representation for stochastic integrals and equations~»,
\emph{SIAM Journal on Control} 4, 1966.
\item I.~Karatzas and S.~E.~Shreve, \emph{Brownian Motion and Stochastic Calculus}, Springer,
2\up{e} éd., 1991.
\end{omsources}

\section{Exercices}

\begin{exercise}[$\star$]\label{exo:qm:ito-calculus:1}
Calculer $\E[(\int_0^1W_t\,dW_t)^2]$ à l'aide de l'isométrie et le vérifier avec l'\cref{ex:qm:ito-calculus:wdw}.
\end{exercise}

\begin{exercise}[$\star$]\label{exo:qm:ito-calculus:2}
Écrire $d(W_t^3)$ à l'aide de la formule d'Itô, puis en déduire $\E[W_t^3]$ et une \omterm{def:qm:probability-at-speed:martingale}{martingale} construite à partir de $W^3$.
\end{exercise}

\begin{exercise}[$\star$]\label{exo:qm:ito-calculus:3}
Une action suit un \omterm{def:qm:brownian-motion:bm}{mouvement brownien} géométrique avec $\mu = 10\%$ et $\sigma = 40\%$. Quelles sont la moyenne
et la médiane de $S_{10}/S_0$, et à quelle volatilité la médiane ne croîtrait pas du tout~?
\end{exercise}

\begin{exercise}[$\star\star$]\label{exo:qm:ito-calculus:4}
Montrer que $e^{W_t - t/2}$ est une \omterm{def:qm:probability-at-speed:martingale}{martingale} et calculer $\E[e^{2W_1}]$.
\end{exercise}

\begin{exercise}[$\star\star$]\label{exo:qm:ito-calculus:5}
Appliquer la règle du produit à $tW_t$ pour écrire $\int_0^1 W_t\,dt$ comme intégrale stochastique, puis calculer
sa variance.
\end{exercise}

\begin{exercise}[$\star\star$]\label{exo:qm:ito-calculus:6}
$X = W^1$ et $Y = 2W^1 + W^2$, avec $W^1, W^2$ indépendants. Calculer $[X, Y]_t$, $[Y]_t$ et la
corrélation instantanée de $X$ et $Y$.
\end{exercise}

\begin{exercise}[$\star\star\star$]\label{exo:qm:ito-calculus:7}
\emph{Programmation.} Refaire le backtest avec une moyenne mobile sur 60 jours~: calculer le ratio de Sharpe
avec look-ahead défini par la \cref{prop:qm:ito-calculus:lookahead}, puis sa moyenne sur 100 trajectoires simulées de dix ans.
\end{exercise}

\begin{exercise}[$\star\star\star$]\label{exo:qm:ito-calculus:8}
\emph{Trouver la faille.} «~Notre règle décide à partir du cours de clôture et négocie
dans l'enchère de clôture~; le backtest lui crédite le rendement du lendemain, donc il n'y a pas de look-ahead.~»
\end{exercise}

\section{Problème~: le look-ahead dans le backtest}

\begin{problem}[{Problème du week-end --- un ratio de Sharpe de 5.4 sur une marche aléatoire}]\label{pb:qm:ito-calculus:1}
Le logarithme du cours est une marche aléatoire de volatilité quotidienne de 1~\%. La règle détient $\theta_t = (S_t -
\bar S^{(20)}_t)/c$, avec $c$ choisi de sorte que $\theta$ soit de variance unitaire~; le backtest erroné crédite $\theta_t$
de l'accroissement $S_t - S_{t-1}$, le backtest honnête de $S_{t+1} - S_t$.

\medskip\noindent\textbf{Partie I --- La règle.}
\begin{enumerate}
\item Quelle est la valeur de $c$~?
\item Quelle est la corrélation $a$ entre $\theta_t$ et l'accroissement du jour~?
\item Quelles sont la moyenne et l'écart-type du gain sur la même barre par jour~?
\item Quel est le ratio de Sharpe annualisé sur la même barre~?
\item Combien le backtest honnête gagne-t-il en moyenne, et pourquoi~?
\end{enumerate}

\medskip\noindent\textbf{Partie II --- La simulation.}
\begin{enumerate}[resume]
\item Quels ratios de Sharpe l'historique de dix ans à graine fixée donne-t-il pour les deux backtests~?
\item Sur 400 historiques, quelles sont la moyenne et l'écart-type de chacun~?
\item Vérifier l'identité~: gain sur la même barre $-$ gain honnête $=$ covariation, sur l'historique à graine fixée.
\item Quelle statistique $t$ la fonction \texttt{lookahead\_test} renvoie-t-elle pour l'historique à graine fixée~?
\item Pourquoi l'écart-type du ratio de Sharpe honnête vaut-il environ $1/\sqrt{10}$~?
\end{enumerate}

\medskip\noindent\textbf{Partie III --- Le point de vue d'Itô.}
\begin{enumerate}[resume]
\item Écrire le P\&L sur la même barre comme une \omterm{def:qm:ito-calculus:integral}{intégrale d'Itô} plus une covariation.
\item Pourquoi l'espérance de la partie d'Itô est-elle nulle~?
\item Quels sont les ratios de Sharpe avec look-ahead pour des fenêtres de 5, 10 et 60 jours~?
\item Si le backtest erroné tournait sur des barres de cinq minutes, soit 78 par jour, avec la même règle appliquée aux barres, quel serait son ratio de Sharpe annualisé~?
\item Négocier à l'ouverture du lendemain plutôt qu'à la clôture supprimerait-il le biais~?
\end{enumerate}

\medskip\noindent\textbf{Partie IV --- Jugement.}
\begin{enumerate}[resume]
\item Comment détecter cette erreur dans le backtest de quelqu'un d'autre~?
\item Le biais de look-ahead est-il toujours positif~?
\item Sur des cours réels comportant un rebond bid-ask, qu'est-ce qui peut encore produire une covariation entre la position et le cours~?
\item Énoncer le \emph{résultat nommé}~: le ratio de Sharpe avec look-ahead de la règle de tendance sur la même barre, sur une marche aléatoire.
\item En une phrase~: qu'est-ce qui fait d'un backtest une \omterm{def:qm:ito-calculus:integral}{intégrale d'Itô}~?
\end{enumerate}
\end{problem}

\section{Questions d'entretien}

\begin{interviewq}[$\star$ \iqroles{chercheur}]\label{iq:qm:ito-calculus:1}
Calculez $\int_0^TW_t\,dW_t$. Pourquoi ne vaut-elle pas $\tfrac12W_T^2$~?
\end{interviewq}

\begin{interviewq}[$\star$ \iqroles{trader, chercheur}]\label{iq:qm:ito-calculus:2}
Une action a un rendement espéré de 10~\% et une volatilité de 40~\%. Quel est son rendement logarithmique espéré,
et que gagne un investisseur type sur dix ans~?
\end{interviewq}

\begin{interviewq}[$\star\star$ \iqroles{chercheur}]\label{iq:qm:ito-calculus:3}
$W_t^2$ est-il une \omterm{def:qm:probability-at-speed:martingale}{martingale}~? Trouvez le processus qu'il faut retrancher, ainsi que la variance de $\int_0^TW\,dW$.
\end{interviewq}

\begin{interviewq}[$\star\star$ \iqroles{développeur, chercheur}]\label{iq:qm:ito-calculus:4}
Comment concevriez-vous un moteur de backtest qui ne puisse pas commettre de look-ahead par construction~?
\end{interviewq}

\begin{interviewq}[$\star\star$ \iqroles{chercheur, banque}]\label{iq:qm:ito-calculus:5}
Qu'est-ce qu'une \omterm{def:qm:ito-calculus:local}{martingale locale} qui n'est pas une \omterm{def:qm:probability-at-speed:martingale}{martingale}, et pourquoi un quant de valorisation devrait-il s'en préoccuper~?
\end{interviewq}

\begin{interviewq}[$\star\star\star$ \iqroles{chercheur, ingénieur ML}]\label{iq:qm:ito-calculus:6}
Itô ou Stratonovich~: lequel la finance utilise-t-elle, lequel la physique utilise-t-elle, et pourquoi~?
\end{interviewq}
